\section{Conclusion}

We started from the copymint detector of Kotzer et al., in which four perceptual hashes vote and two
agreeing votes mean a copy, and we set out to make it faster. Reimplementing it showed that one
of the four signals, sHash, compares images one way and breaks the triangle inequality, so the
BK-tree that searches it can silently skip a real copy.

We replaced sHash with ORB feature matching and changed nothing else. On the geometric edits
sHash was built for, ORB scores F1 90.6 against sHash's 76.4, even with sHash given a
threshold chosen on the test set; at equal precision it catches 90.0\% of copies against
27.2\%. In the full detector at the paper's own thresholds, the swap raises F1 from
77.5 to 86.4 at the same 98.1\% precision, catching about 1,500 more of 11,400 copies. The
gain holds when the paper's thresholds are re-tuned (+5.8), on images with two edits (+8.7),
and modestly on the authors' own images (+3.3).

The improvement has a measured price. ORB needs about 14.5\,KB per image where the paper's
four hashes need about 56 bytes, and no descriptor budget lets it both fit in a mint
transaction and beat sHash. Using ORB means changing where features are stored.

Our edit classifier predicts accurately which signals an edit breaks, yet it does not improve
the detector, because broken signals go silent rather than voting wrongly. The one situation
where it could help, detail destroyed and colour changed together, is not observable with our
cues because pixelation erases the colour trace.

In short: we replaced one signal, measured it fairly, and reported its cost.
